
    %%%%%%%%%%%%%%%%%%%%%%%%%%%%%%%%%%%%%%%%%%%%%%%%%%%%%%%%%%%%%%%%%%%%%%%%%%%%%%%%
    %%% Classe do documento
    \documentclass[12pt, addpoints]{exam}
    \UseRawInputEncoding

    %%%%%%%%%%%%%%%%%%%%%%%%%%%%%%%%%%%%%%%%%%%%%%%%%%%%%%%%%%%%%%%%%%%%%%%%%%%%%%%%
    % preambulo.tex
    %
    % Arquivo de configuração de packages para uso o LaTeX, conforme minhas
    % preferências e modelos pessoais.
    %
    % Para maiores informações, visite:
    %    https://github.com/abrantesasf/latex
    %
    % NÃO ALTERE SE NÃO SOUBER O QUE ESTÁ FAZENDO!
    %%%%%%%%%%%%%%%%%%%%%%%%%%%%%%%%%%%%%%%%%%%%%%%%%%%%%%%%%%%%%%%%%%%%%%%%%%%%%%%%


    %%%%%%%%%%%%%%%%%%%%%%%%%%%%%%%%%%%%%%%%%%%%%%%%%%%%%%%%%%%%%%%%%%%%%%%%%%%%%%%%
    %%% Estruturas de controle:
    \input{static/utils/controles}

    %%%%%%%%%%%%%%%%%%%%%%%%%%%%%%%%%%%%%%%%%%%%%%%%%%%%%%%%%%%%%%%%%%%%%%%%%%%%%%%%
    %%% Compilação condicional em PDF:
    \input{static/utils/pdf}

    %%%%%%%%%%%%%%%%%%%%%%%%%%%%%%%%%%%%%%%%%%%%%%%%%%%%%%%%%%%%%%%%%%%%%%%%%%%%%%%%
    %%% Configurações de layout da página:
    \input{static/utils/layout}

    %%%%%%%%%%%%%%%%%%%%%%%%%%%%%%%%%%%%%%%%%%%%%%%%%%%%%%%%%%%%%%%%%%%%%%%%%%%%%%%%
    %%% Configurações para autores:
    \input{static/utils/autores}

    %%%%%%%%%%%%%%%%%%%%%%%%%%%%%%%%%%%%%%%%%%%%%%%%%%%%%%%%%%%%%%%%%%%%%%%%%%%%%%%%
    %%% Cabeçalho e rodapé:
    %%% Atenção! É MUITO DIFÍCIL ajustar apropriadamente os running headers
    %%% e running footers da primeira vez. Você pode confiar no LaTeX e deixar que
    %%% ele faça o ajuste automaticamente ou pode quebrar a cabeça e
    %%% tentar melhorar algo que o LaTeX já faz muito bem, mesmo que não fique
    %%% exatamente do jeito que você quer. O que você prefere? Se quiser ajustar
    %%% manualmente, descomente a linha abaixo e faça as configurações no arquivo
    %%% "utils/headings.tex".
    %\input{static/utils/headings}

    %%%%%%%%%%%%%%%%%%%%%%%%%%%%%%%%%%%%%%%%%%%%%%%%%%%%%%%%%%%%%%%%%%%%%%%%%%%%%%%%
    %%% Configuração para as fontes do documento:
    %%% Se você quiser usar fontes próprias ou do sistema, precisará ajustar as
    %%% configurações no arquivo "utils/fontes.tex".
    %%% ATENÇÃO: por padrão eu utilizo XeLaTeX com fontes proprietárias projetadas
    %%% por Matthew Butterick, adquiridas em https://mbtype.com, que não podem ser
    %%% distribuídas devido à licença de utilização. Se você usar XeLaTeX, você
    %%% DEVERÁ OBRIGATORIAMENTE ajustar as fontes no arquivo "utils/fontes.tex".
    \input{static/utils/fontes}

    %%%%%%%%%%%%%%%%%%%%%%%%%%%%%%%%%%%%%%%%%%%%%%%%%%%%%%%%%%%%%%%%%%%%%%%%%%%%%%%%
    %%% Sumário:
    \input{static/utils/toc}

    %%%%%%%%%%%%%%%%%%%%%%%%%%%%%%%%%%%%%%%%%%%%%%%%%%%%%%%%%%%%%%%%%%%%%%%%%%%%%%%%
    %%% Matemática:
    \input{static/utils/matematica}

    %%%%%%%%%%%%%%%%%%%%%%%%%%%%%%%%%%%%%%%%%%%%%%%%%%%%%%%%%%%%%%%%%%%%%%%%%%%%%%%%
    %%% Notas de rodapé:
    \input{static/utils/footnotes}

    %%%%%%%%%%%%%%%%%%%%%%%%%%%%%%%%%%%%%%%%%%%%%%%%%%%%%%%%%%%%%%%%%%%%%%%%%%%%%%%%
    %%% Referências cruzadas, links, citações:
    \input{static/utils/crossrefs}

    %%%%%%%%%%%%%%%%%%%%%%%%%%%%%%%%%%%%%%%%%%%%%%%%%%%%%%%%%%%%%%%%%%%%%%%%%%%%%%%%
    %%% Configurações para os hiperlinks em PDF:
    \input{static/utils/pdfconfig}

    %%%%%%%%%%%%%%%%%%%%%%%%%%%%%%%%%%%%%%%%%%%%%%%%%%%%%%%%%%%%%%%%%%%%%%%%%%%%%%%%
    %%% Glossário e índice remissivo:
    \input{static/utils/glossindex}

    %%%%%%%%%%%%%%%%%%%%%%%%%%%%%%%%%%%%%%%%%%%%%%%%%%%%%%%%%%%%%%%%%%%%%%%%%%%%%%%%
    %%% Referências bibliográficas:
    \input{static/utils/bibliografia}

    %%%%%%%%%%%%%%%%%%%%%%%%%%%%%%%%%%%%%%%%%%%%%%%%%%%%%%%%%%%%%%%%%%%%%%%%%%%%%%%%
    %%% Cores:
    \input{static/utils/cores}

    %%%%%%%%%%%%%%%%%%%%%%%%%%%%%%%%%%%%%%%%%%%%%%%%%%%%%%%%%%%%%%%%%%%%%%%%%%%%%%%%
    %%% Computação:
    \input{static/utils/computacao}

    %%%%%%%%%%%%%%%%%%%%%%%%%%%%%%%%%%%%%%%%%%%%%%%%%%%%%%%%%%%%%%%%%%%%%%%%%%%%%%%%
    %%% Gráficos:
    \input{static/utils/graficos}

    %%%%%%%%%%%%%%%%%%%%%%%%%%%%%%%%%%%%%%%%%%%%%%%%%%%%%%%%%%%%%%%%%%%%%%%%%%%%%%%%
    %%% Ícones:
    \input{static/utils/icones}

    %%%%%%%%%%%%%%%%%%%%%%%%%%%%%%%%%%%%%%%%%%%%%%%%%%%%%%%%%%%%%%%%%%%%%%%%%%%%%%%%
    %%% Blocos:
    \input{static/utils/blocos}

    %%%%%%%%%%%%%%%%%%%%%%%%%%%%%%%%%%%%%%%%%%%%%%%%%%%%%%%%%%%%%%%%%%%%%%%%%%%%%%%%
    %%% Boxes coloridos:
    \input{static/utils/colorbox}

    %%%%%%%%%%%%%%%%%%%%%%%%%%%%%%%%%%%%%%%%%%%%%%%%%%%%%%%%%%%%%%%%%%%%%%%%%%%%%%%%
    %%% Tabelas:
    \input{static/utils/tabelas}

    %%%%%%%%%%%%%%%%%%%%%%%%%%%%%%%%%%%%%%%%%%%%%%%%%%%%%%%%%%%%%%%%%%%%%%%%%%%%%%%%
    %%% Ambientes de listas:
    \input{static/utils/listas}

    %%%%%%%%%%%%%%%%%%%%%%%%%%%%%%%%%%%%%%%%%%%%%%%%%%%%%%%%%%%%%%%%%%%%%%%%%%%%%%%%
    %%% Ambientes floats e similares:
    \input{static/utils/floats}

    %%%%%%%%%%%%%%%%%%%%%%%%%%%%%%%%%%%%%%%%%%%%%%%%%%%%%%%%%%%%%%%%%%%%%%%%%%%%%%%%
    %%% Ajustes para a classe EXAM:
    \input{static/utils/exam}

    %%%%%%%%%%%%%%%%%%%%%%%%%%%%%%%%%%%%%%%%%%%%%%%%%%%%%%%%%%%%%%%%%%%%%%%%%%%%%%%%
    %%% Ajustes para textos:
    \input{static/utils/texto}

    %%%%%%%%%%%%%%%%%%%%%%%%%%%%%%%%%%%%%%%%%%%%%%%%%%%%%%%%%%%%%%%%%%%%%%%%%%%%%%%%
    %%% Ambientes, aliases, comandos e customizações em geral para serem
    %%% utilizadas nos documentos. Defina aqui qualquer customização específica
    %%% para seus documentos!
    \input{static/utils/customizacoes}

    %%%%%%%%%%%%%%%%%%%%%%%%%%%%%%%%%%%%%%%%%%%%%%%%%%%%%%%%%%%%%%%%%%%%%%%%%%%%%%%%
    %%% Ajuste do layout, espaçamento de linhas e etc.:
    \geometry{a4paper, portrait, left=2.0cm, right=2.0cm, top=2.0cm, bottom=2.0cm}
    \singlespacing

    %%%%%%%%%%%%%%%%%%%%%%%%%%%%%%%%%%%%%%%%%%%%%%%%%%%%%%%%%%%%%%%%%%%%%%%%%%%%%%%%
    %%% Configurações de cabeçalho e rodapé (não use fancyhdr pois ocorre conflito):
    \pagestyle{headandfoot}
    \runningheadrule
    \firstpageheader{}{}{}
    \runningheader{Sem disciplina}{}{Avaliação}
    \firstpagefooter{}{}{Página \thepage\ de \numpages}
    \runningfooter{}{}{Página \thepage\ de \numpages}
    \runningfootrule

    %%%%%%%%%%%%%%%%%%%%%%%%%%%%%%%%%%%%%%%%%%%%%%%%%%%%%%%%%%%%%%%%%%%%%%%%%%%%%%%%
    %%% Configurações para as propriedades do PDF:
    \hypersetup{
    hidelinks,           % Comente para web, descomente para impressão
    colorlinks=true,      % True para web, False para impressão
    pdftitle=llsls: Avaliação,
    pdfauthor=brenno,
    pdfsubject=Sem disciplina,
    pdfkeywords=['Sem tag'],
    pdfinfo={
        CreationDate={}, % Ex.: D:AAAAMMDDHH24MISS
        ModDate={}       % Ex.: D:AAAAMMDDHH24MISS
    }
    }
    \input{static/utils/pdfconfig}

    %%%%%%%%%%%%%%%%%%%%%%%%%%%%%%%%%%%%%%%%%%%%%%%%%%%%%%%%%%%%%%%%%%%%%%%%%%%%%%%%
    %%% Começa o documento
    \begin{document}

    %%%%%%%%%%%%%%%%%%%%%%%%%%%%%%%%%%%%%%%%%%%%%%%%%%%%%%%%%%%%%%%%%%%%%%%%%%%%%%%%
    %%% Folha de instruções padronizadas da UVV para a prova
    \pdfbookmark[1]{Instruções}{instrucoes}
    \begin{figure}[H]
        \begin{center}
            \includegraphics[scale=0.3]{{static/imgs/logo-uvv.png}}
        \end{center}	
    \end{figure}

    \vspace{-1cm}

    \begin{table}[H]
        \centering
        \begin{tabular}{|llll|}
            \hline
            \multicolumn{2}{|l|}{\textbf{Disciplina:} Sem disciplina} & 
            \multicolumn{1}{l|}{\multirow{3}{*}{\textbf{\makecell{Nota:\\ \,\\ \,}}}} & 
            \multirow{3}{*}{\textbf{\makecell{Coordenador:\\ \,\\ \,}}} \\ 
            \cline{1-2}
            \multicolumn{2}{|l|}{\textbf{Professor:} brenno \hspace{4.72cm} } & 
            \multicolumn{1}{l|}{} & \\ 
            \cline{1-2}
            \multicolumn{2}{|l|}{\textbf{Aluno:}} & 
            \multicolumn{1}{l|}{} & \\ 
            \hline
            \multicolumn{1}{|l|}{\textbf{Turma:} \hspace{6.3cm} } & 
            \multicolumn{1}{l|}{\textbf{Semestre:}} & 
            \multicolumn{2}{l|}{\textbf{Valor:} 10 pontos} \\ 
            \hline
            \multicolumn{1}{|l|}{\textbf{Data:}} & 
            \multicolumn{3}{l|}{\textbf{Avaliação:}} \\ 
            \hline
        \end{tabular}
    \end{table}

    \begin{center}
        \fbox{\setlength\fboxsep{0.3cm} \fbox{\parbox{15.4cm}{
            \begin{center}
                \textbf{INSTRUÇÕES PARA A REALIZAÇÃO DESTA AVALIAÇÃO:}
            \end{center}
            \begin{itemize}[leftmargin=*]
                \item Esta é a primeira avaliação da disciplina Sem disciplina, 
                    e engloba o conteúdo referente Sem tag.
                \item Esta avaliação contém 10 questões objetivas, \textbf{todas obrigatórias}.
                \item \textbf{Leia atentamente cada questão!} As questões objetivas terão uma,
                    e apenas uma única, resposta a ser marcada.
                \item \textbf{Desligue o celular} antes de começar. A avaliação é
                    \textbf{individual} e \textbf{sem consulta}.
                \item As questões podem ser respondidas, na prova, com lápis ou caneta. Ao final
                    da prova \textbf{{VOCÊ DEVERÁ PREENCHER O GABARITO COM CANETA
                    DE COR PRETA}} \textbf{\underline{OU AZUL}} para que seja feita a correção
                    automatizada através da leitura do padrão de respostas marcado no
                    gabarito.
                \item \textbf{CUIDADO ao preencher o gabarito} pois eles são \textbf{INDIVIDUAIS
                    e NOMEADOS} (cada estudante tem um gabarito próprio específico, com um
                    código de identificação único). \textbf{Questões rasuradas são anuladas}
                    pelo software de leitura do gabarito.
                \item Ao final da prova, \textbf{DEVOLVA A PROVA E O GABARITO} para o professor.
                \item Siga todas as normas de \textbf{Integridade Acadêmica} da disciplina.
                    Alunos que forem flagrados com qualquer espécie de ``cola'' ou trocando
                    informações com outros alunos terão suas avaliações recolhidas, as notas
                    zeradas, e a situação será encaminhada para a coordenação para a aplicação
                    das penalidades previstas pela universidade.
                \item Com 90 minutos de avaliação você terá tempo de até 9,min por questão,
                    em média.
                \item Boa avaliação!
            \end{itemize}
            \vspace{2cm}
        }}}
    \end{center}
    \newpage

    %%%%%%%%%%%%%%%%%%%%%%%%%%%%%%%%%%%%%%%%%%%%%%%%%%%%%%%%%%%%%%%%%%%%%%%%%%%%%%%%
    %%% Inicia o bloco de questões:
    \begin{questions}

    \newpage
    \question
Ao executar o seguinte comando SQL, nenhum dado foi retornado, ou seja, não
ocorreu nenhum erro de sintaxe mas nenhum dado foi encontrado:
\begin{tcolorbox}
 \begin{verbatim}1   SELECT c.nome AS cliente
 2   FROM clientes c
 3   LEFT JOIN pedidos p ON (p.cliente_id = c.cliente_id)
 4   WHERE p.cliente_id IS NULL
 5   ORDER BY cliente;
 \end{verbatim}
\end{tcolorbox}
Por que esse comando SQL não retornou nenhum dado?

\begin{choices}
\choice A junção que deveria ter sido realizada, na linha 3, era uma RIGHT JOIN ao invés de uma LEFT JOIN.
\choice A junção que deveria ter sido realizada, na linha 3, era uma FULL JOIN ao invés de uma RIGHT JOIN.
\choice Há um erro semântico na linha 3
\choice Todos os clientes cadastrados fizeram, pelo menos, um pedido de compra.
\choice Há um erro lógico na linha 4.
\end{choices}

\question
Considere as afirmações abaixo sobre endereços IP:
\begin{itemize}
    \item[I.] Uma rede IP classe C fornece até 256 endereços válidos para serem atribuídos a equipe.
    \item[II.] A quantidade máxima de bits que pode ser utilizada para se definir sub-redes em uma rede IP classe C é seis (6).
    \item[III.] A máscara padrão para uma rede classe B é 255.255.255.0.
\end{itemize}
Qual das seguintes alternativas representa as afirmações corretas?
\begin{choices}
\choice Somente I e II.
\choice Somente II.
\choice Somente III.
\choice Somente I.
\choice Somente II e III.
\end{choices}

\question
Redes semânticas, frames e lógica são formalismos utilizados principalmente em: 
\begin{choices}
\choice redes neurais
\choice descoberta de conhecimento em bases de dados
\choice inferência em sistemas especialistas
\choice representação de conhecimento
\choice IA distribuída
\end{choices}

\question
Um agente SNMP é um aplicativo que é executado:
\begin{choices}
\choice em roteadores com filtragem de pacotes com o objetivo de proteger acesso a rede
\choice em computadores denominados de gerentes
\choice em um dispositivo de rede
\choice em "firewalls" com o objetivo de proteger acesso a rede
\choice a partir de um computador específico para monitorar a rede
\end{choices}

\question
Considere o algoritmo da busca sequencial de um elemento em um conjunto com \( n \) elementos. A expressão que representa o tempo médio de execução desse algoritmo para uma busca bem sucedida é:
\begin{choices}
\choice \( n(n + 1)/2 \)
\choice \( (n + 1)/2 \)
\choice \( n^2 \)
\choice \( n \log n \)
\choice \( \log_2 n \)
\end{choices}

\question
Considere um problema em que são dados 5 objetos com os seguintes pesos e valores:

\[
\text{pesos: } (W_1, W_2, W_3, W_4, W_5) = (6, 10, 9, 5, 12)
\]
\[
\text{valores: } (P_1, P_2, P_3, P_4, P_5) = (8, 5, 10, 15, 7)
\]

Além disso, é dada uma mochila que suporta até 30 unidades de peso, para transportar os objetos. O objetivo do problema é preencher a mochila de tal forma que o valor total dos objetos a serem transportados seja o maior possível, mas sem exceder o limite de peso suportado pela mochila. Assuma que é permitido colocar fração de um objeto na mochila. Qual das seguintes alternativas corresponde ao valor máximo obtido no preenchimento da mochila:
\begin{choices}
\choice 30.34
\choice 21.5
\choice 12.2
\choice 43.1
\choice 38.83
\end{choices}

\question
O processo de visualização de objetos 3D envolve uma série de passos desde a representação vetorial de um objeto até a exibição da imagem correspondente na tela do computador (pipeline 3D). Selecione a alternativa abaixo que reflete a ordem correta em que esses passos devem ocorrer. 
\begin{choices}
\choice Recorte 3D, transformação de câmera, rasterização, projeção, mapeamento para coordenadas de tela. 
\choice Transformação de câmera, mapeamento para coordenadas de tela, recorte 3D, rasterização, projeção. 
\choice Projeção, transformação de câmera, recorte 3D, mapeamento para coordenadas de tela, rasterização. 
\choice Nenhuma das respostas acima está correta.
\choice Transformação de câmera, recorte 3D, projeção, mapeamento para coordenadas de tela, rasterização. 
\end{choices}

\question
Uma característica de uma arquitetura RISC é:
\begin{choices}
\choice A arquitetura é constituída de milhares de processadores
\choice O processador é formado por válvulas e transistores
\choice Uma arquitetura de alto risco pois o mercado de hardware evolui muito rapidamente
\choice Possui um pequeno conjunto de instruções simples
\choice Possui um grande conjunto de instruções complexas
\end{choices}

\question
(Adaptado de ENADE 2011) Considere o diagrama abaixo, que ilustra um banco de
dados que descreve os empregados que trabalham nos departamentos de uma empresa:
\begin{figure}[H]
  \centering
  \caption{Empregados e departamentos}
  \vspace{-0.2cm}
  \label{fig:empregados}
  %\fbox{
    \includegraphics[width=0.5\linewidth]{static/imgs/defaul.png}
  %}\\
  %\footnotesize{Fonte: xxxx.}
\end{figure}
\vspace{-0.4cm}
Na linguagem SQL, o comando mais simples para recuperar os códigos dos
departamentos que contém empregados que ganham menos do que 2000 e mais do que
5000 reais é:

\begin{choices}
\choice \begin{verbatim}SELECT cod_departamento
FROM empregados
WHERE salario < 2000 AND salario > 5000;
\end{verbatim}

\choice \begin{verbatim}SELECT cod_departametno
FROM empregados
WHERE salario <= 2000 OR salario >= 5000;
\end{verbatim}

\choice \begin{verbatim}SELECT cod_departamento
FROM empregados
WHERE salario <= 2000 AND salario >= 5000;
\end{verbatim}

\choice \begin{verbatim}SELECT cod_departamento
FROM empregados
WHERE salario < 2000 OR salario > 5000;
\end{verbatim}

\choice \begin{verbatim}SELECT cod_departamento
FROM empregados
WHERE salario BETWEEN 2000 AND 5000;
\end{verbatim}

\end{choices}

\question
Considere as seguintes afirmações sobre resolução de problemas em IA.
\begin{enumerate}[label=\Roman*.]
    \item A* é um conhecido algoritmo de busca heurística.
    \item O Minimax é um dos principais algoritmos para jogos de dois jogadores, como o xadrez.
    \item Busca em espaço de estados é uma das formas de resolução de problemas em IA.
\end{enumerate}
São corretas:
\begin{choices}
\choice Apenas I e III
\choice Apenas II e III
\choice Apenas I e II
\choice I, II e III
\choice Apenas III
\end{choices}



    %%%%%%%%%%%%%%%%%%%%%%%%%%%%%%%%%%%%%%%%%%%%%%%%%%%%%%%%%%%%%%%%%%%%%%%%%%%%%%%%
    %%% Finaliza o bloco de questões:
    \end{questions}
    \end{document}
    